\chapter{Conclusiones}

Esta práctica ha supuesto un cambio de registro respecto a las anteriores: en lugar de configurar una herramienta o reproducir un ataque controlado en laboratorio, el ejercicio consiste en \textit{navegar el ecosistema real de divulgación de vulnerabilidades}, aprender a leerlo con criterio y posicionarse técnicamente frente a fallos publicados por terceros. Esa diferencia es, creo, el aprendizaje principal que me llevo: en seguridad operativa, una parte muy importante del trabajo no es generar conocimiento nuevo sino \textit{consumir bien} el que generan otros (MITRE, NIST, fabricantes, ZDI, CISA), priorizarlo con criterio y traducirlo a acciones concretas sobre el parque propio.

Lo que más me ha sorprendido al analizar las tres CVE elegidas es comprobar que un mismo \textit{Base Score} de 9.8 con un vector idéntico esconde realidades muy distintas en cuanto a riesgo real. Sobre el papel, las tres son intercambiables: ataque remoto, sin privilegios, sin interacción, con impacto alto en CIA. En la práctica, sin embargo, una afecta a una librería gráfica omnipresente en cualquier Windows del mundo, otra a un parser RTP de una app de mensajería ya parcheada hace años, y la tercera a un agente de gestión de SAI presente sólo en CPDs concretos. El CVSS dice \textit{lo mismo} de las tres y la respuesta racional de un responsable de seguridad debe ser distinta para cada una. Esto resuena con algo que ya intuía: el CVSS Base es un \textit{punto de partida} útil para filtrar, pero la priorización fina depende de variables que no están en el vector ---exposición del activo, criticidad del negocio, existencia de exploit funcional, presencia en KEV--- y que sólo se obtienen cruzando varias fuentes.

Otra observación relevante es la importancia de \textit{seguir periódicamente} las publicaciones de NVD, INCIBE-CERT y CISA. La frecuencia con la que aparecen vulnerabilidades críticas en componentes nucleares ---ya sea GDI+ en Windows, parsers de protocolos en mensajería, agentes de monitorización en Linux--- hace inviable enterarse \textit{a posteriori} de los problemas relevantes; la única respuesta razonable es la suscripción sistemática (boletines, RSS de NVD, listas oficiales del fabricante) y la inclusión del CVE-management como tarea continua del equipo de operaciones. INCIBE-CERT, en particular, es un punto de entrada útil en el contexto español porque traduce, contextualiza y filtra ese flujo para el tejido productivo nacional, y CISA aporta la perspectiva de \textit{explotación real} con el KEV.

Las dificultades han venido sobre todo de la \textbf{limitación del buscador del NIST}, ya mencionada en el capítulo 2: filtros que no admiten rangos mayores de 120 días han obligado a trocear la consulta en ventanas de tres meses, repitiéndolas dos veces por ventana (una para \textit{Published Date} y otra para \textit{Last Modified}). En la práctica esto implicó del orden de ocho ejecuciones del buscador para cubrir un único año natural y un trabajo manual de seguimiento de las ventanas ya recorridas. Una segunda dificultad ha sido distinguir entre \textit{informe primario} (CVE.org), \textit{informe enriquecido} (NVD) y \textit{aviso del fabricante}: aunque los tres apuntan a la misma vulnerabilidad, los datos no siempre coinciden con exactitud (estado del registro, fecha de parcheo, severidad asignada por el CNA frente a la del NIST), por lo que ha sido necesario aprender a contrastar siempre al menos dos fuentes antes de dar por cerrado cualquier dato concreto.

En conjunto, la práctica refuerza una idea sencilla pero importante: el ciclo CVE/CVSS no es burocracia, sino \textit{infraestructura compartida} de la comunidad de seguridad. Saber moverse en él ---identificar fuentes oficiales, leer vectores con criterio, distinguir lo crítico de lo cosmético, priorizar con datos--- es un componente básico del oficio de cualquier responsable de seguridad de la información.
